% Generated draft: requires mathematical and source review.
\documentclass[11pt]{article}
\usepackage[margin=25mm]{geometry}
\usepackage{fontspec}
\setmainfont{Latin Modern Roman}
\usepackage{amsmath,amssymb,mathtools,unicode-math}
\setmathfont{Latin Modern Math}
\usepackage{enumitem,needspace,xurl,hyperref}
\setlength{\parindent}{0pt}
\setlength{\parskip}{5pt}
\setlength{\emergencystretch}{3em}
\newcommand{\N}{\mathbb{N}}
\newcommand{\Z}{\mathbb{Z}}
\newcommand{\Q}{\mathbb{Q}}
\newcommand{\R}{\mathbb{R}}
\newcommand{\C}{\mathbb{C}}
\newcommand{\F}{\mathbb{F}}
\newcommand{\A}{\mathbb{A}}
\newcommand{\PP}{\mathbb{P}}
\newcommand{\E}{\mathbb{E}}
\newcommand{\eps}{\varepsilon}
\newcommand{\dd}{\,\mathrm{d}}
\begin{document}
\section*{UnsolvedMath ID: 20002703}
\textbf{Source problem code:} catalogue entry
\subsection{Definitions and mathematical statement}
For an irreducible finite real reflection group $W$ of rank $n$, Coxeter number $h$, and exponents $e_i$ (invariant degrees $d_i=e_i+1$), the source defines
\[
 \operatorname{Cat}^{(k)}(W)=\prod_{i=1}^n\frac{kh+e_i+1}{e_i+1}
 =\prod_i\frac{kh+d_i}{d_i},\qquad k\in\Z_{>0}.
\]
For reducible groups the componentwise product convention is needed. The source asks for natural Fuss generalizations of $NC_W,NN_W,\Delta_W$, structural explanations of their enumerative refinements, an exponent/invariant-theoretic framework, and analogues of its Problems 1.1 and 1.2. It describes Armstrong's graded join-semilattice of $k$-divisible noncrossing partitions and its Fuss--Narayana ranks, Fomin--Reading generalized-cluster facets and $h$-numbers, type-$A,B$ $(k+2)$-angulation models, Athanasiadis's affine-arrangement regions and co-filtered root-ideal multichains for Weyl groups, Haiman's coroot-quotient orbit count, and Sommers's Lie-theoretic occurrences. The remarks ask for Fuss $M,F,H$ triangles and applications to Garside structures, cluster algebras, and free probability; their final application sentence is truncated.

The aid treats only $W=A_1^n$, every $k\ge1$ and $n\ge0$, with total count $(k+1)^n$. A convention-explicit Boolean state model chooses $a_i\in\{0,1,\ldots,k\}$. Put $s_i$ into the unique delta factor $\delta_{a_i}$, so $\delta_0\cdots\delta_k$ is the Coxeter element and total reflection length is $n$. Associate the cluster facet having $-\alpha_i$ if $a_i=0$, and the positive root $\alpha_i$ of color $a_i$ otherwise. Associate a decreasing chain of filters by including $\alpha_i$ at every level if $a_i=0$, and at levels $r<a_i$ otherwise. This is a bijective threshold parametrization in the Boolean root poset. Its final filter contains exactly the roots with $a_i=0$, and
\[
 \operatorname{Fix}(\delta_0)
 =\bigcap_{\alpha\text{ in final filter}}\alpha^\perp
 =\bigcap_{-\alpha_i\text{ in facet}}\alpha_i^\perp.
\]
A colored enumerator records each state separately as $\prod_i(z_{i,0}+\cdots+z_{i,k})$; weighting state $0$ by $z_i$ and the other states by $1$ gives $\prod_i(k+z_i)$.

The supplied factorizations, with $F$ counting positive and negative face vertices, are
\[
 F(x,y)=(1+kx+y)^n,\qquad H(x,y)=(k+xy)^n,\qquad
 M(x,y)=(k-ky+xy)^n.
\]
Here $H$ uses final-filter/rank-$k$ indecomposable counts, which coincide coordinatewise in the Boolean case. For $M$ use the product of rank-one posets with $k$ minimal rank-zero states and one maximal rank-one state, and the convention $\sum_{u\le v}\mu(u,v)x^{\operatorname{rk}u}y^{\operatorname{rk}v}$. Each coordinate contributes $k-ky+xy$.

Rank conventions require care: delta-zero reflection rank has polynomial $(k+t)^n$, whereas the usual cluster $h$-polynomial from this $F$ is $(1+kt)^n$. A shelling taking the negative vertex first in each rank-one factor has restriction size equal to the number of positive vertices, hence $n-\ell_T(\delta_0)$. Thus these statistics agree after rank reversal, not literally under the displayed delta-zero rank. The aid advertises a statistic-preserving shelling interpretation; its full conventions must be checked before identifying that phrase with an unreversed rank equality.

For $k>1$ the rank-one standard Fuss poset has several minimal elements and no meet for two distinct ones. It is not itself a lattice, obstructing a verbatim identification with a Garside simple-element lattice. This does not exclude modified posets or other Garside applications; the $n=0$ singleton case is exceptional. The Boolean construction does not provide the requested uniform general-type framework.
\subsection{Short English statement}
The source seeks a general explanation of Fuss--Catalan families and applications. The aid gives a colored Boolean model with exact flats and triangle factorizations. Its rank-to-shelling comparison needs an explicit reversal convention, and its rank-one nonlattice obstruction concerns only a verbatim Garside-simple interpretation.
\subsection{Sources}
\begin{itemize}
\item Local UnsolvedMath registry, permanent ID 20002703, source code catalogue entry; created and updated 14 August 2026. Canonical \texttt{aim-probability-notes.json}, zero-based \texttt{notes[144]}, tag \texttt{problem}, item 1.3.
\item AIM ``Braid groups, clusters and free probability'': \url{https://aimath.org/WWN/braidgroups/braidgroups.pdf}. Statement status is reconstructed\_unverified; supplied verification reports restoration of notation/formula and notes a stray page 6 and truncated remarks. Original literature field is empty.
\item Original references: Armstrong [1]; Athanasiadis [2,3]; Fomin--Reading [21]; classical type-$A$ formulas [21,30]; Tzanaki [51], also named for generalized $H$ triangles; Haiman [28]; Sommers [46]. Full bibliography and the text of referenced Problem 1.2 are absent.
\item Research attempt 1: \texttt{partial\_result}, problem status at run \texttt{partially\_solved}, rigor label \texttt{proved}, novelty confidence low. Ulam AI / UnsolvedMath Contributors, CC BY 4.0; machine-generated aid, not peer reviewed, independently verify before citation.
\end{itemize}
\end{document}
